\section{Introduction}
Large-scale language models (LLMs) show excellent performance on various tasks~\citep{gpt3,opt}. However, serving LLMs is budget and energy-consuming due to their gigantic model size. For example, the GPT-3~\citep{gpt3} model contains 175B parameters, which will consume at least 350GB of memory to store and run in FP16, requiring 8$\times$48GB A6000 GPUs or 5$\times$80GB A100 GPUs just for inference.
Due to the huge computation and communication overhead, the inference latency may also be unacceptable to real-world applications.
\emph{Quantization} is a promising way to reduce the cost of LLMs~\cite{dettmers2022llmint8, zeroquant}. By quantizing the \emph{weights and activations} with low-bit integers, we can reduce GPU memory requirements, in size and bandwidth, and accelerate compute-intensive operations (\ie, \texttt{GEMM}\footnote{General matrix multiply} in linear layers, \texttt{BMM}\footnote{Batch matrix multiply} in attention). For instance, INT8 quantization of weights and activations can halve the GPU memory usage and nearly double the throughput of matrix multiplications compared to FP16.

\begin{figure}
    \centering
    \includegraphics[width=0.9\linewidth]{figures/trend.pdf}
    \caption{The model size of large language models is developing at a faster pace than the GPU memory in recent years, leading to a big gap between the supply and demand for memory. Quantization and model compression techniques can help bridge the gap. }
    \vspace{-1em}
    \label{fig:trend}
\end{figure}


However, unlike CNN models or smaller transformer models like BERT~\cite{bert}, the \emph{activations} of LLMs are difficult to quantize. When we scale up LLMs beyond 6.7B parameters, systematic outliers with large magnitude will emerge in activations~\cite{dettmers2022llmint8}, leading to large quantization errors and accuracy degradation.
ZeroQuant~\citep{zeroquant} applies dynamic per-token activation quantization and group-wise weight quantization (defined in \fig{fig:quantize_notation} \sect{sec:quantization}). It can be implemented efficiently and delivers good accuracy for GPT-3-350M and GPT-J-6B. However,  it can not maintain the accuracy for the large OPT model with 175 billion parameters (see Section~\ref{sec:accurate-quantization}).
\texttt{LLM.int8()}~\cite{dettmers2022llmint8} addresses that accuracy issue by further introducing a mixed-precision decomposition (\ie, it keeps outliers in FP16 and uses INT8 for the other activations). However, it is hard to implement the decomposition efficiently on hardware accelerators.
Therefore, deriving an \emph{efficient}, \emph{hardware-friendly}, and preferably \emph{training-free} quantization scheme for LLMs that would use INT8 for all the compute-intensive operations remains an open challenge.

We propose \method, an accurate and efficient post-training quantization (PTQ) solution for LLMs.
\method relies on a key observation: even if activations are much harder to quantize than weights due to the presence of outliers~\cite{dettmers2022llmint8},  different tokens exhibit similar variations across their channels.
\begin{figure}
    \centering
    \includegraphics[width=0.88\linewidth]{figures/intuition.pdf}
    \caption{\method's intuition: the activation $\mathbf{X}$ is hard to quantize because outliers stretch the quantization range, leaving few effective bits for most values. We migrate the scale variance from activations to weights $\mathbf{W}$ during offline to reduce the quantization difficulty of activations. The smoothed activation $\hat{\mathbf{X}}$ and the adjusted weight  $\hat{\mathbf{W}}$ are both easy to quantize. 
       }
    \label{fig:intuition}
    \vspace{-10pt}
\end{figure}

Based on this observation, \method offline migrates the quantization difficulty from activations to weights (\fig{fig:intuition}).
\method proposes a mathematically equivalent per-channel scaling transformation that significantly smooths the magnitude across the channels, making the model quantization-friendly.
Since \method is compatible with various quantization schemes, we implement three efficiency levels of 
quantization settings for \method (see \tbl{tab:settings}, O1-O3).
Experiments show that \method is hardware-efficient: it can maintain the performance of OPT-175B~\cite{opt}, BLOOM-176B~\cite{scao2022bloom} , GLM-130B~\cite{zeng2022glm}, and MT-NLG 530B~\cite{smith2022using},
leading to up to 1.51$\times$ speed up and 1.96$\times$ memory saving on PyTorch.
\method is easy to implement. We integrate \method into FasterTransformer, the state-of-the-art transformer serving framework, achieving up to 1.56$\times$ speedup and halving the memory usage compared with FP16.
Remarkably, \method allows serving large models like OPT-175B using only half number of GPUs compared to FP16 while being faster, and enabling the serving of a 530B model within one 8-GPU node.
Our work democratizes the use of LLMs by offering a turnkey solution to reduce the serving cost. We hope \method can inspire greater use of LLMs in the future.
